\documentclass[11pt]{article}
\usepackage[margin=1in]{geometry}
\usepackage[T1]{fontenc}
\usepackage{lmodern}
\usepackage{amsmath,amssymb,amsthm}
\usepackage{hyperref}

\newtheorem{theorem}{Theorem}
\newtheorem{lemma}{Lemma}

\newcommand{\Pettis}{\mathcal P}
\newcommand{\ND}{\mathcal{ND}}
\newcommand{\cL}{\mathcal L}
\newcommand{\eps}{\varepsilon}

\title{A Hilbert-Subspace Spaceability Result for Non-Differentiable Pettis Primitives}
\author{Run fa\_banach\_001, agent\_super\_00}
\date{2026-07-01}

\begin{document}
\maketitle

\section*{Source Problem}

Bongiorno--Darji--Di Piazza \cite{BDD} prove that, for every infinite
dimensional Banach space \(X\), the set \(\ND_X\) of strongly measurable
Pettis integrable functions \(f\colon[0,1]\to X\) whose Pettis primitive is
nowhere weakly differentiable is lineable.  They ask whether \(\ND_X\) is
spaceable.

We use the standard Pettis norm
\[
  \|f\|_{\Pettis}
  =
  \sup_{E\in\cL}
  \left\|(P)\int_E f\,d\lambda\right\|,
\]
with functions identified in the usual Pettis-norm sense.  The source paper
explicitly discusses this norm topology and notes that it is not complete in
general.

\section*{Partial Result}

\begin{theorem}
Let \(X\) be a Banach space containing a subspace isomorphic to \(\ell_2\).
Then \(\ND_X\) is spaceable in the Pettis norm topology.
\end{theorem}

\section*{Idea of the Proof}

The source paper's \(\ell_2\)-case uses dyadic intervals, pairwise disjoint
positive-measure sets \(A(\sigma,i)\subset I_\sigma\), and an orthonormal
family \(e(\sigma,i)\).  Its lineability argument chooses many branches whose
coordinate choices are eventually almost disjoint.

For closedness, we instead put the \(n\)th generator permanently in the
coordinate column \(i=n\), beginning at dyadic level \(n\).  A factor
\(\sqrt{n+1}\) exactly compensates for the tail
\(\sum_{k\ge n}(k+1)^{-2}\).  This makes the Pettis norm of the generated
function equivalent to the \(\ell_2\)-norm of the coefficient sequence, hence
the generated subspace is closed.  The first nonzero coefficient still appears
on every sufficiently fine dyadic scale, forcing unbounded primitive
difference quotients at every point.

\section*{Proof}

We first prove the result for \(X=\ell_2\).  Use the notation from the source
paper's Section 3.  Thus \(2^{<\infty}\) is the dyadic tree, \(I_\sigma\) is
the dyadic interval associated to \(\sigma\), and, for every pair
\((\sigma,i)\) with \(0\le i\le |\sigma|\), \(A(\sigma,i)\subset I_\sigma\)
are pairwise disjoint measurable sets of positive measure.  Let
\(\{e(\sigma,i)\}\) be the corresponding orthonormal basis of \(\ell_2\).

For \(n\ge0\), define
\[
g_n
=
\sqrt{n+1}
\sum_{k=n}^{\infty}
\sum_{\sigma\in 2^k}
 \frac{1}{(k+1)2^{k/2}}\,
 \frac{1}{\lambda(A(\sigma,n))}\,
 e(\sigma,n)\chi_{A(\sigma,n)} .
\]
For \(a=(a_n)\in\ell_2\), set
\[
G_a=\sum_{n=0}^{\infty} a_n g_n.
\]
Equivalently,
\[
G_a=
\sum_{k=0}^{\infty}\sum_{\sigma\in2^k}\sum_{n=0}^{k}
a_n\sqrt{n+1}\,
\frac{1}{(k+1)2^{k/2}}\,
\frac{1}{\lambda(A(\sigma,n))}
e(\sigma,n)\chi_{A(\sigma,n)} .
\]

For a measurable set \(E\), write
\[
\theta_{\sigma,n}(E)=
\frac{\lambda(E\cap A(\sigma,n))}{\lambda(A(\sigma,n))},
\qquad 0\le \theta_{\sigma,n}(E)\le1.
\]
Then
\[
(P)\int_E G_a
=
\sum_{k=0}^{\infty}\sum_{\sigma\in2^k}\sum_{n=0}^{k}
a_n\sqrt{n+1}\,
\frac{\theta_{\sigma,n}(E)}{(k+1)2^{k/2}}\,
e(\sigma,n).
\]
By orthogonality,
\[
\left\|(P)\int_E G_a\right\|_2^2
\le
\sum_{k=0}^{\infty}\sum_{\sigma\in2^k}\sum_{n=0}^{k}
\frac{|a_n|^2(n+1)}{(k+1)^2 2^k}
=
\sum_{n=0}^{\infty}|a_n|^2(n+1)
\sum_{k=n}^{\infty}\frac1{(k+1)^2}.
\]
For every \(n\ge0\),
\[
1
\le
(n+1)\sum_{k=n}^{\infty}\frac1{(k+1)^2}
\le
2.
\]
Hence
\[
\|G_a\|_{\Pettis}\le \sqrt2\,\|a\|_2.
\]
Taking \(E=[0,1]\) gives equality in the coefficient factors
\(\theta_{\sigma,n}(E)=1\), so the same estimate also yields
\[
\|G_a\|_{\Pettis}\ge \left\|(P)\int_{[0,1]}G_a\right\|_2
\ge \|a\|_2.
\]
Thus \(T\colon a\mapsto G_a\) is an isomorphic embedding of \(\ell_2\) into
the Pettis-normed space \(\Pettis([0,1],\ell_2)\).  Since \(\ell_2\) is
complete and the induced norm is equivalent to \(\|\cdot\|_2\), \(T(\ell_2)\)
is closed in the Pettis norm.

It remains to check that every nonzero \(G_a\) belongs to \(\ND_{\ell_2}\).
Let \(a\ne0\), and let \(n_0\) be the least index with \(a_{n_0}\ne0\).  Let
\(F_a(t)=(P)\int_0^t G_a\).  Fix \(x\in[0,1]\) and \(M>0\).  For every
sufficiently small \(h>0\), choose a dyadic interval \(I_\tau\subset[x,x+h]\)
of length \(2^{-j}\), with \(j\ge n_0\), as in the source proof; then
\(h<4\cdot 2^{-j}\) and \(j\to\infty\) as \(h\downarrow0\).  Since
\(I_\tau\subset[x,x+h]\) and the coordinates are orthogonal,
\[
\left\|F_a(x+h)-F_a(x)\right\|_2
\ge
\left\|(P)\int_{I_\tau}G_a\right\|_2.
\]
The \(n_0\)-column alone gives
\[
\left\|(P)\int_{I_\tau}G_a\right\|_2
\ge
|a_{n_0}|\sqrt{n_0+1}\,
2^{-j/2}
\left(\sum_{m=j}^{\infty}\frac1{(m+1)^2}\right)^{1/2}.
\]
Therefore
\[
\left\|\frac{F_a(x+h)-F_a(x)}{h}\right\|_2
\ge
\frac{|a_{n_0}|\sqrt{n_0+1}}{4}\,
2^{j/2}
\left(\sum_{m=j}^{\infty}\frac1{(m+1)^2}\right)^{1/2}.
\]
The right-hand side tends to \(+\infty\), since the tail sum is at least a
constant multiple of \(1/(j+1)\).  The same argument applies for \(h<0\).
Thus the primitive has unbounded norm difference quotients at every point.
Weak differentiability at a point would make the corresponding difference
quotients weakly convergent, hence norm bounded by uniform boundedness, a
contradiction.  So \(G_a\in\ND_{\ell_2}\).

Now let \(X\) contain a subspace \(Y\) isomorphic to \(\ell_2\), and let
\(U\colon\ell_2\to Y\subset X\) be an isomorphism.  The map
\[
a\in\ell_2\longmapsto U\circ G_a
\]
is still an isomorphic embedding in the Pettis norm, because
\[
\|U^{-1}\|^{-1}\|G_a\|_{\Pettis}
\le
\|U G_a\|_{\Pettis}
\le
\|U\|\|G_a\|_{\Pettis}.
\]
Bounded linear images of Pettis integrable functions are Pettis integrable,
and the unbounded norm-difference-quotient estimate survives with the lower
constant \(\|U^{-1}\|^{-1}\).  Hence every nonzero element of this closed
infinite-dimensional subspace belongs to \(\ND_X\).  This proves the theorem.

\section*{Verification Notes}

The construction uses only the dyadic sets and orthogonal-coordinate
calculation already present in the source paper.  The new point is the
column-wise choice of generators and the \(\sqrt{n+1}\) normalization, which
turns the Pettis norm estimates into a two-sided \(\ell_2\) estimate.

The result is partial.  It covers Hilbert spaces and all Banach spaces
containing an isomorphic copy of \(\ell_2\).  It does not settle the source
problem for arbitrary infinite-dimensional Banach spaces.

\section*{Novelty Check}

Cheap run indexes were searched for arXiv:1403.1908, title terms, Pettis
primitive phrases, and spaceability terms.  Web searches for exact title and
phrases such as ``Pettis primitives spaceable'' and ``nowhere weakly
differentiable Pettis spaceability'' found only the source paper.  A local
full-source search found citations to the source paper and general
Pettis-integration references, but no later exact solution of the spaceability
problem.

\section*{Human Review Recommendation}

Review as a likely valid partial result.  The main points to check are the
Pettis-norm closedness argument and the transfer from \(\ell_2\) to an
isomorphic Hilbert subspace of \(X\).

\begin{thebibliography}{9}
\bibitem{BDD}
B.~Bongiorno, U.~B. Darji, and L.~Di Piazza,
\emph{Lineability of non-differentiable Pettis primitives},
arXiv:1403.1908.

\bibitem{DU}
J.~Diestel and J.~J. Uhl,
\emph{Vector Measures},
Mathematical Surveys, vol.~15, American Mathematical Society, 1977.
\end{thebibliography}

\end{document}
